\section{网络模型与验证方法}

\subsection{信号通道与测量平面}
\label{sec:model-plane}

每条信号通道为 $N$ 级梯形网络，由正值串联电阻和正值对地电容组成，并由理想电压源驱动。图~\ref{fig:circuit} 展示二级实例，输出取末端节点电压差。传递函数按零初始条件定义。缓冲后的比值模型中，样品电阻先改变输入源幅度，滤波元件保持固定；输入前史确定相容的初始状态。

差模输入和共模输入定义为

\begin{equation}\label{eq:1}
v_d=v_p-v_n,\qquad v_c=\frac{v_p+v_n}{2}.
\end{equation}

在差模输入为零时，共模转差模增益为 \(G_{DC}\) = \(v_{d,\mathrm{out}}/v_{c,\mathrm{in}}\)；在共模输入为零时，差模增益为 \(G_{DD}\) = \(v_{d,\mathrm{out}}/v_{d,\mathrm{in}}\)。共模抑制比为 \(\lvert G_{DD}/G_{DC}\rvert\)。对于传递函数分别为 \(H_P\) 和 \(H_N\) 的独立通道，

\begin{equation}\label{eq:2}
G_{DC}=H_P-H_N,\qquad
G_{DD}=\frac{H_P+H_N}{2}.
\end{equation}

共模抵消要求 \(G_{DC}\) 作为频率的有理函数恒等于零。

\begin{figure}[htbp]
\centering
\includegraphics[width=0.86\linewidth]{figures/Fig_01_circuit.pdf}
\caption{两条对地二阶 RC 支路及输出端接。每条梯形支路含两个串联电阻和两个对地电容。可选电容 $C_d$ 跨接输出，跨接路径以蓝色点线表示。两个输入由对地理想电压源驱动。}
\label{fig:circuit}
\end{figure}


\subsection{系数表示与负载条件}
\label{sec:coefficients}

以 \(R_*\) 和 \(C_*\) 为基准归一化元件值，定义

\begin{equation}\label{eq:3}
r_i=R_i/R_*,\qquad q_i=C_i/C_*,\qquad x=sR_*C_*.
\end{equation}

对于二阶示例，支路的四个系数为

\begin{equation}\label{eq:4}
\begin{aligned}
a&=r_1q_1+(r_1+r_2)q_2,\\
b&=r_1r_2q_1q_2,\\
c&=r_1+r_2,\\
d&=r_1r_2q_1.
\end{aligned}
\end{equation}

由两个节点方程可得开路电压传递函数 \(H_0\) 和输出阻抗 \(Z_0\)：

\begin{equation}\label{eq:5}
H_0(x)=\frac{1}{1+ax+bx^2},\qquad
Z_0(x)=R_*\frac{c+dx}{1+ax+bx^2}.
\end{equation}

输出端与地之间接入电阻 \(R_L\) 时，其归一化电导为 \(g=R_*/R_L\)，带载传递函数为

\begin{equation}\label{eq:6}
H_g(x)=\frac{H_0(x)}{1+Z_0(x)/R_L}
      =\frac{1}{1+gc+(a+gd)x+bx^2}.
\end{equation}

每一种负载条件下，两条通道均连接相同的指定电阻 \(R_L\)。统计匹配支路 N 的可行元件组合时，参考支路 P 保持固定。每个电导条件下比较的对象均为完整传递函数；比较两个电导条件，即分别比较每条通道在这两种负载下随频率变化的传递函数。

\subsection{精确运算与数值验证}

采用精确有理数运算核查系数恒等式和示例元件组合。电路仿真使用 LTspice 26.1.1，取 \(R_*=1\,\mathrm{k\Omega}\)、\(C_*=10\,\mathrm{nF}\)。扫频范围为 10 Hz 至 10 MHz，每十倍频程 24 点，共 145 点。共模激励在两个输入端同时施加 1 V；差模激励分别施加 +0.5 V 和 -0.5 V。

37 个仿真算例覆盖一阶对照、二阶支路、不同元件匹配策略、有限源电阻、对称与非对称电阻端接，以及输出跨接电容。根据保存的网表独立建立节点方程，将所得复节点电压与仿真输出逐点比较。电阻负载分类的解析结论采用第~\ref{sec:model-plane}、\ref{sec:coefficients} 节所述的理想电压源和两侧等值电阻负载假设。


一般阶矩阵包括 $N=1,2,3,4,6,8,12$ 的84例原生交流仿真。独立节点核查采用实际元件值及仿真器规定的元件模型；其中，未显式给出串联电阻的电感在 LTspice 中默认串入1~m$\Omega$。另做两组成对诊断运行，分别显式指定零或1~m$\Omega$。瞬态验证采用第~\ref{sec:dynamic-verif} 节所列的波形与热身窗口，两个四级算例的细化检查将最大步长减半。

\FloatBarrier
